\documentclass[10pt]{article}

\usepackage[utf8]{inputenc}

\usepackage[T1]{fontenc}

\usepackage{amsmath}

\usepackage{amsfonts}

\usepackage{amssymb}

\usepackage[version=4]{mhchem}

\usepackage{stmaryrd}

\usepackage{mathrsfs}



\begin{document}

тодов, дающих оценки для количества почти простых чисел, не превосходящих $x$ и принадлежащих последовательности натуральных чисел $\left\{a_{n}\right\}$. Методы решета нозволяют получать такие оденки, если известны "хорошие» приближения в среднем для количества чисел $a_{n} \leqslant x, \quad$ принадлежащих прогрессиям, модуль к-рых растет вместе с $x$. Среди методов решета, развитых после работ В. Бруна, особое значение приобрело Сельберга решето. Наиболее сильные результаты получают сочетанием методов решета с аналитическими. Метод решета в сочетании с IIIнирельмана методом дал возможность әффективно найти значение $k$ такое, что любое натуральное число $n \geqslant 4$ можно представить суммой нө более $k$ шростых чисел.



Два целых числа $a$ и $b$ наз. сравнимыми по модулю $m \geqslant 1$, если они дают одинаковые остатки при делении на число $m$. Для өтого отношения на множестве целых чисел К. Гаусс $(1801)$ ввел запись $a \equiv b(\bmod m)$. Әта форма записи, выявляющая аналогию сравнений с уравнениями, оказалась удобной и способствовала развитию сравнений теории.



Многие результаты, полученные до этого П. Ферма (P. Fermat), Л. Эйлером (L. Euler), Ж. Лагранжем (J. Lagrange) и др., а также китайская теорема об осmaткаx, просто формулируются и доказываются на языке теории сравнений. Один из наиболее интересных результатов этой теории - квадратичный закон взаи.нности.



Древним вавилонянам было известно большое количество «пифарагоровых троек», видимо, им был известен какой-то способ, позволяющий находить сколько угодно целочисленных решений неопределенного уравнения $x^{2}+y^{2}=z^{2}$. Пифагорейцы пользовались формулами $x=$ $=\frac{a^{2}-1}{2}, y=a, z=\frac{a^{2}+1}{2}$ для отыскания таких решениі̆ этого уравнения. Евклид указал метод, позволягщий последовательно находить целочисленные решения уравнения $x^{2}-2 y^{2}=1$ (частный случай Пелля уравнения). В «Арифметике» Диофанта (3 в. н. э.) заметна попытка построения теории неопределенных уравнений (см. Диобантовь уравнения). В час'ности, Диофант чри решении уравнений 2-й степени и нек-рых уравнений высших степеней спстематически пользуется приемом, к-рый позволяет ему находить из одного рационального решения данного уравнения др. рационалыьные решения того же уравнения. П. Ферма (17в.) открыл новый метод - «метод спуска» и решил ряд уравнений таким методом, но объявленная им как решенная Ферма теорема великая, оказалась не под силу элементарным методам.



П. Ферма решил задачу о шредставлении натуральных чисел суммой двух квадратов целых чисел. В результате исследований Ж. Јагранжа (1773) и К. Гаусса (1801) была решена задаqа о представлении чисел определенвой бинарной көадратичной формой. К. Гаусс построил общую теорию бинарных квадратичных форм. Решение задачи о представлении чисел формами более высокоӥ степени (напр., Варинга проблема) и квадратичными формами с большим числом переменных обычно выходит за рамки әлементарных методов. Только нек-рые частные случаи таких задач решаются әлементарно. Примером может служить т е о р е м а ЈI а г р а н ж а: каждое натуральное число есть сумма четырех квадратов целых чиселг. Следует заметить, что в своей «Арифметике» Диофант неоднократно пользовался возможностью представить натуральное число суммой четырех квадратов целых чисел.



К Ә. т. ч. относят теорию разбиений, основы к-рой были заложены JI. Әйлером (1751). Одна из основных задач теории разбиений - изучение функции $P(n)$ числа представлений натурального числа $n$ в виде сум- мы натуральных слагаөмых. В теории разбиений рассматривают и др. функции подобного типа.



Депные дроби появились в связи с задачами приближенных вычислений (извлечение квадратного корня из натурального числа, отыскание приближения действительных чисел обыкновенными дробями с малыми знаменателями). Цепные дроби находят применение при решении неопределенных уравнений 1-й и 2-й степени, Используя апшарат цешных дробей, И. Ламберт (J. Lambert, 1766) впервые установил иррациональность числа л. При решении различных задач о приближении действительных чисел рациональными числами в Ә. т. ч. наряду с цепными дробями используют также Дирихле принцип.



В теории чисел легко назвать большое число әлементарно формулируемых и нерешенных до сих пор задач. Напр.: конечно или нет множество четных совершенных чисел; существует ли хотя бы одно нечетное совершенное число; конечно или нет множество простых чисел Ферма; конечно или бесконечно множество простых Мерсенна чисел; конечно или нет множество простых чисел вида $n^{2}+1 ;$ верно ли, что между квадратами любых соседних натуральных чисел содержится хотя бы одно простое число; ограничено или нет множество неполных частных разложения $\sqrt[3]{2}$ в цепную дробь.



Jum.: [1] В е н к о в Б. А., Элементарная теория чиселt, М.- Л., 1937; [2] В ин о г р а д о в И. М., Основы теории чисел, юизд., Млементарные методы в аналитической теории чисел, М., $1962 ;$ [4] Х и н ч и н А. Я., Цепные дроби, 4 изд.. М., 1978; [5] Г а у с с К. Ф., Труды по теории чисел, пер. с

лат., М., 1959; [6] Д ә в е н п о р т Г., Высшая арифметика, пер.' с англ., М., 1965 ; [7] Т р о с т Ә., Простые числа, пер. нем., М., 1959; [8] Э н д p ю с $\Gamma$., Теория разбиений, пер. с ангјі., М., 1982; [9] История математики, т. 1-3, М., 1970-72; [10] В и л е й т н е р Г., История математики от Декарта до середины XIX столетия, пер. c нем., M., $1966 ;[11] \mathrm{D}$ i c k-

s on L. E., History of the theory of numbers, V. $1-3, \mathrm{~N}$. Y., 1971; [12] H a rd y G. H., W $\mathbf{r}$ i g $\mathbf{h} \mathbf{t} \mathbf{E}$. M., An introduction to the theory of numbers, 5 ed., Oxf., 1979; [13] S i e r p in s$\mathrm{k}$ i W., Elementary theory of numbers, Warsz., 1964.



A. A. Бухитаб,' B. И. Нечаев.



ӘЈЕМЕНТАРНОЕ СОБЫТИЕ - исходное ионя'ие вероятностной модели. В определении вероя'ностного пространства $(\Omega, \mathscr{A}, \mathrm{P})$ непустое множество $\Omega$ наз. пространством Э. с., а его любая точка $\omega \in \Omega$ наз. э ли ем ен т а р н ы м с о б ы т е м. При неформальном подходе множество $\Omega$ описывает множество всех исходов нек-рого случайного әксперимента и Ә. с. $\omega$ соответствуе'і элементарному исходу: эксперимент заканчивается одним и только одним элементарным исходом, эти исходы неразложимы и взаимно исключаю' друг друга. Существует принцищиальная разница мељду $\%$. с. $\omega-$ точкой множества $\Omega$ и событием $\{\omega\}$ - элементом некрого класса множеств $\mathscr{A}$. См. Вероятностей теория, Вероятностное пространство, Случайное событиче.



A. B. IIpoxipos.



ЭЛЕМЕНТАРНЫЕ ДЕЈИТЕЛИ м а т р и ц Ы $F(x)$ над колыцом многочленов $k[x]$ - стешени унитарных неприводимых многочленов над полем $k$, на к-рые разлагаются инвариантные множители матрицы $F(x)$. Две $(m \times n)$-матрицы над $k[x]$, имеющие один и тот же ранг, тогда и только тогда әквивалентны (т. е. нолучаются одна из другої с помощью әлементарных операций), котда они обладают одной и той же системой Ә. Д.



$\ni$ лг м е т т а р ны ми д е ли те л я ми $(n \times n)-$ матрицы $A$ над полем $k$ наз. Э. Д. ее характеристич. матрицы $\left\|x E_{n}-A\right\|$. Они могут быть получены следующим образом. Пусть $D_{l}(x)$ - наибольшиї общий делитель миноров порядка $l$ матрицы $x E_{n}-A, 1 \leqslant l \leqslant n$ и $D_{0}=1$. Тогда инвариантными множителями матрицы $\left\|x E_{n}-A\right\|$ служа'



\[

i_{l}(x)=\frac{D_{l}(x)}{D_{l-1}(x)}, \quad i=1, \ldots ; n .

\]





\end{document}